% ============================================================================
% Doc. 269 -- FFGFT and RA/PMT: a bidirectional structural comparison
%
% Compares FFGFT with the RA/PMT architectural scheme (Guevara 2026) in
% both directions: (A) FFGFT evaluates RA/PMT using the six-axis
% comparison operator from Doc. 265; (B) the RA/PMT grid is applied to
% FFGFT. Honest caveat: the two frameworks are not on the same level.
%
% Author: Johann Pascher
% Date: June 2026
% Status: Comparison document. Related: Doc. 265 (comparison operator),
%         Doc. 267 (cosmological degeneracy), Doc. 203 (recursion R).
% ============================================================================

\section*{Purpose}

This document compares FFGFT (Fundamental Fractal Geometric Field Theory,
T$^4$ time-mass duality) with the RA/PMT scheme of J.~T.~Guevara
(\emph{Regional Arenas / Perception Mechanism Theory}, 2026, three parts:
consciousness, fundamental physics, cosmology).

The comparison uses the same explicit comparison operator as Doc.~265 --
but applied here \emph{in both directions}:

\begin{itemize}
  \item \textbf{Direction A:} FFGFT evaluates RA/PMT using the
        six-axis operator (Doc.~265).
  \item \textbf{Direction B:} the RA/PMT grid itself
        (RA4 $\to$ RA1.5 $\to$ RA2 $\to$ RA1) is applied to FFGFT --
        just as Guevara applies it to QFT, GR and emergent spacetime.
\end{itemize}

Applying it both ways is exactly the point of an exposed operator
(Doc.~265): each side may render the other with its own tool, and both
renderings remain examinable side by side.

\textbf{Up front -- the most important caveat:} FFGFT and RA/PMT are
\emph{not on the same level}. FFGFT is a physical theory with one free
parameter ($\xi = 1/7500$) from which concrete numbers follow (CMB peak
ratios, $H_0$, mass relations). RA/PMT is, in Guevara's own words,
\emph{``not a new physical theory''} -- a meta-scheme for classifying
theories. The comparison is therefore not between two competing models
of the same system (as Doc.~265 was for FFGFT/Blume-Maleev), but between
a theory and a theory-ordering scheme. This shapes every finding below.

% ============================================================================
\section{The RA/PMT scheme in brief}
% ============================================================================

RA/PMT decomposes every theory into five layers:

\begin{center}
\begin{tabular}{p{2.2cm}p{9.5cm}}
\toprule
\textbf{Layer} & \textbf{Function} \\
\midrule
RA4 & Blueprint -- the structured space of all \emph{possible}
      configurations (prior to any realization). \\
\addlinespace
RA3 & Residual -- the theory's undefined terms, explicitly declared as
      unresolved rather than hidden. \\
\addlinespace
RA2 & Mapping -- the rules/equations that unfold a chosen blueprint into
      physical appearance. \\
\addlinespace
RA1.5 & Mechanism -- the \emph{selector} that picks one configuration
      from RA4. (HPM for consciousness, PSM for physics, CSM for
      cosmology -- the same role, per Guevara.) \\
\addlinespace
RA1 & Expression -- the physically realized. \\
\bottomrule
\end{tabular}
\end{center}

RA/PMT's central thesis: QM, GR, emergent spacetime and cosmology all
possess RA4 and RA2 but \emph{no} RA1.5. The absence of this selector is
claimed to be the common root of the measurement problem, the
initial-conditions problem, fine-tuning, the arrow of time, and the
cosmological constant. RA/PMT ``completes'' the theories by adding RA1.5 --
without altering their equations.

% ============================================================================
\section{Direction A: FFGFT evaluates RA/PMT (six-axis operator)}
% ============================================================================

We apply the comparison operator from Doc.~265: two frameworks are
compared along six axes (shared invariants, operators, mathematical
structures, predictions, physical mechanisms, hardware implications).
On each axis: shared, partial, or not.

\begin{table}[htbp]
\centering
\caption{FFGFT $\leftrightarrow$ RA/PMT along the six comparison axes.
$\checkmark$ = shared, $\sim$ = only at the meta-level, $\times$ = not shared.}
\label{tab:directionA}
\begin{tabular}{p{6.5cm}c p{4cm}}
\toprule
\textbf{Comparison axis} & \textbf{Status} & \textbf{Note} \\
\midrule
1. Invariants          & $\times$ &
  RA/PMT supplies no invariant; FFGFT: $\tilde{T}\cdot m=1$, winding number. \\
2. Operators           & $\sim$ &
  Structural only: RA/PMT's ``selector'' vs.\ FFGFT's recursion
  $\mathcal{R}$ (Doc.~203) -- but $\mathcal{R}$ is defined, RA1.5 is not. \\
3. Mathematical structure & $\sim$ &
  Shared only as \emph{layering discipline}: FFGFT's cores $>$ bridges
  $>$ reductions $\approx$ RA/PMT's RA4 $>$ RA2 $>$ RA1. \\
4. Predictions (quantitative) & $\times$ &
  RA/PMT produces not a single number; FFGFT: peak ratios, $H_0$,
  masses from $\xi$. \\
5. Physical mechanisms & $\times$ &
  RA1.5/CSM is pre-physical and unspecified; FFGFT mechanisms are
  geometrically concrete. \\
6. Hardware implications & $\times$ &
  RA/PMT makes no statement about devices/measurements. \\
\bottomrule
\end{tabular}
\end{table}

\textbf{Finding, Direction A:} On four of six axes there is \emph{no}
contact, on two only contact at the meta-level (axes~2 and~3). By the
classification of Doc.~265 this is not even Level~3 (complementary in the
usual sense) -- it is \textbf{category difference}: the two frameworks do
not answer the same kind of question. The only genuine overlap is
methodological: both separate cleanly what follows rigorously from what
remains open. FFGFT does this through its three-layer methodology (cores
proved from $\xi$ $>$ bridges proved algebraically $>$ reductions as
plausibility), RA/PMT through RA4/RA3/RA2/RA1.5/RA1. This discipline is
congenial and related -- but it is a \emph{stance}, not shared physics.

% ============================================================================
\section{Direction B: the RA/PMT grid applied to FFGFT}
% ============================================================================

Guevara applies his grid to other theories himself (QFT, GR, emergent
spacetime, IIT, \ldots). The logical consequence: one can apply it to
FFGFT too. This is illuminating.

\begin{center}
\begin{tabular}{p{2.2cm}p{9.5cm}}
\toprule
\textbf{Layer} & \textbf{FFGFT assignment} \\
\midrule
RA4 & The T$^4$ mode space: winding numbers $\vec{n}\in\mathbb{Z}^4$,
      the possible modes. \\
\addlinespace
RA3 & Open points, already explicitly declared in FFGFT
      (P-register, Doc.~190): e.g.\ the exponent $41/4$ (R20),
      the rigorous $C_\ell$ projection (R30). \\
\addlinespace
RA2 & The derivations from $\xi$: Jacobi degeneracy, Bose-Einstein,
      3D projection, recursion operator $\mathcal{R}$ (Doc.~203). \\
\addlinespace
RA1.5 & \textbf{Here RA/PMT asks: where is your selector?} \\
\addlinespace
RA1 & Observable: CMB peak ratios, $H_0$, mass relations. \\
\bottomrule
\end{tabular}
\end{center}

\textbf{The decisive point -- FFGFT has no RA1.5, and needs none.}
Guevara's grid diagnoses a ``missing selector'' in every theory. For
FFGFT the answer is not ``the selector is missing'' but \emph{``the
premise does not hold''}:

\begin{itemize}
  \item RA1.5 presupposes that there is a space of \emph{many possible}
        universes/configurations (RA4) from which \emph{one} is selected.
        That is the multiverse-like assumption RA/PMT carries tacitly
        (as $\Lambda$CDM carries its initial conditions).
  \item FFGFT is \textbf{static}: no Big Bang, no $t=0$, no initial state
        to be selected. The fractal T$^4$ torus is not \emph{one selected}
        among many configurations -- it is \emph{the} structure. There is
        no blueprint space for a selector to choose from.
  \item Where Guevara's cosmology (Part~III) reads the Big Bang,
        inflation and the low initial entropy as ``the expression of a
        selected blueprint'', FFGFT says: these events are not initial
        conditions to be selected, because there is no beginning
        (Doc.~267, inverted initial condition).
\end{itemize}

\textbf{What Direction~B exposes:} The RA/PMT grid is not neutral. It
carries a hidden assumption -- that \emph{every} theory needs a selector
between possibility and actuality. FFGFT is precisely the counterexample:
a theory that deliberately lacks this layer because it does not make the
possibility/actuality split. The ``deficiency'' the grid diagnoses is,
from the FFGFT side, not a deficiency but the absence of a superfluous
assumption.

This is methodologically the most valuable finding of the whole
comparison: it shows that RA1.5 itself -- by Guevara's own RA3
definition -- is a \emph{residual}. ``CSM selects the low-entropy
configuration'' renames the puzzle instead of solving it: there is no
equation, no selection criterion, no prediction. A named selector
without operationalization meets exactly the RA3 condition (``lacks a
defined mechanism'') that RA/PMT charges against other theories.

% ============================================================================
\section{Synthesis: two opposite poles of the same question}
% ============================================================================

Taken together, the two directions give a clear picture. FFGFT and
RA/PMT are neither competitors nor variants -- they are \textbf{opposite
poles of the same basic question}: how does the observed reality come
about?

\begin{center}
\begin{tabular}{p{4cm}p{4cm}p{4cm}}
\toprule
 & \textbf{RA/PMT} & \textbf{FFGFT} \\
\midrule
Basic assumption & many possible universes (RA4), one is selected &
  one static fractal T$^4$, no selection \\
\addlinespace
Selector RA1.5 & necessary but undefined (CSM/PSM/HPM) &
  absent, not needed \\
\addlinespace
Big Bang, inflation, initial entropy & real events whose configuration
  is selected & do not exist ($t=0$ drops out) \\
\addlinespace
$\Lambda$, fine-tuning & selected parameters &
  premise rejected (no dark sector; redshift =
  fractal path-lengthening, Doc.~190 R14a) \\
\addlinespace
Delivers numbers? & no (meta-scheme) & yes ($\xi \to$ peaks, $H_0$, masses) \\
\addlinespace
Strength & layering discipline, conceptual hygiene &
  parameter-free, falsifiable predictions \\
\bottomrule
\end{tabular}
\end{center}

\noindent Where RA/PMT posits a missing layer, FFGFT says: the layer is
not missing, it is not needed -- the premise (many possible
configurations, one selected) is itself a $\Lambda$CDM-like assumption.
This is the same methodological line as the ``inverted initial
condition'' and the symmetric circularity of Doc.~267: RA/PMT's CSM is
essentially $\Lambda$CDM's initial conditions with a metaphysical
selector in front; FFGFT removes the need for both.

\textbf{Fair appraisal.} This is not meant to disparage RA/PMT. As a
\emph{meta-scheme} -- as a discipline for sorting a theory's terms
cleanly into blueprint, residual, mapping, mechanism and expression --
it is a useful diagnostic tool, and its honesty in declaring residuals
deserves credit. It shares with FFGFT the stance ``inspectable rather
than persuasive'' (Doc.~265). The disagreement is not in the method but
in the premise: whether reality is a selection from the possible
(RA/PMT) or a single, static, fractal structure with no selection
(FFGFT).

\textbf{Residuals of this comparison (honestly declared).}
\begin{itemize}
  \item The comparison treats RA/PMT as a meta-scheme. Should Guevara
        later operationalize RA1.5 (give an equation, a selection
        criterion), Direction~A would need re-evaluation -- there might
        then be axis-4 contact.
  \item Direction~B assumes FFGFT's static premise is correct. That is
        an FFGFT-internal position, not a proven theorem; an RA/PMT
        proponent would read it as itself a selected configuration. Both
        readings remain side by side (cf.\ the provenance principle,
        Doc.~265).
\end{itemize}

\noindent\textbf{Related:} Doc.~265 (comparison operator, six axes),
Doc.~267 (cosmological degeneracy, inverted initial condition),
Doc.~203 (recursion operator $\mathcal{R}$), Doc.~190 (R14a, R20, R30).
RA/PMT: J.~T.~Guevara, \emph{Architectural Completion} (Parts~I--III),
June 2026.
